\begin{document}

{
\setcounter{tocdepth}{2}
\tableofcontents
\clearpage
}
\section{TFPT / Universalraum: den verlorenen Spin-Lift wieder
einsetzen}\label{tfpt-universalraum-den-verlorenen-spin-lift-wieder-einsetzen}

\subsection{Forschungsfortsetzung v1.6.10 - 15. September
2026}\label{forschungsfortsetzung-v1.6.10---15.-september-2026}

\textbf{Ergebnis:} Ein bereits in TFPT vorhandenes Umlaufvorzeichen
verbindet drei bisher getrennte Fragen: Welche Information verlieren
Paarmessungen? Welche Fermionrichtungen sind wirklich gekoppelt? Welcher
Grundzustand wird bevorzugt? Für einen ausdrücklich festgelegten
Viererzyklus mit dem unveränderten inneren Tensor W wird diese
Verbindung hier bis zu einem eindeutigen Grundzustand des vollständigen
Hamiltonoperators und seiner isolierten geladenen Antwort verfolgt.

\textbf{Geltungsbereich:} Die Quellenanordnung und ihre Identifikation
mit einer physischen TFPT-Seam sind nicht hergeleitet. Der neue
Grundzustandssatz gilt bei \(0<|g|/\Delta\le1/2000\), der präzisere
Antwortsatz bei \(0<|g|/\Delta\le1/10000\). Das ist \textbf{nicht} der
bisherige Prüfpunkt \(g/\Delta=1/20\). Es gibt keinen neuen
\(\mu N\)-Term. Die Bosonbesetzung wird nicht abgeschnitten. Kein
vollständiges T1-T8-Gate ist geschlossen.

Die Beweise sind analytische Forschungsresultate mit exakten endlichen
Kontrollen und rationalen Fehlerschranken, kein
Proof-Assistant-Zertifikat und kein unabhängiges Fachgutachten. Die alte
lokale 64-Fermion-Bank und die neue globale Vier-Bank-Anordnung sind
verschiedene Hamiltonverträge.

\textbf{Verschärfte Auswahlprüfung:} Abschnitt 15 zeigt, dass Clock und
eindeutiger Grundzustand allein das Umlaufvorzeichen nicht wählen. Ein
vorgegebener gemeinsamer Links-rechts-Reflexionsvertrag würde dagegen
innerhalb der reellen Zweinachbar-Klasse die gleiche Gewichtung
erzwingen; mit dem Ausschluss freier Fermionrichtungen bleibt dann der
negative Lift. Der entsprechende ursprüngliche TFPT-Operatorvertrag ist
noch nicht nachgewiesen.

\subsection{1. Was die erneute Grundlagenlektüre
verändert}\label{was-die-erneute-grundlagenlektuxfcre-veruxe4ndert}

Die erneute Lektüre umfasst die Grundkonstruktion in
\codeword{README.md}, \codeword{docs/THEORY.md},
\codeword{origin\_theory.tex}, die D5/A3- und Pascal-Konstruktion, die
aktuelle Standardmodell-Masterformel und ihre Projektionsgrenzen, den
Kaskadenvergleich im E8-Audit, die Frontier-Verträge sowie die
einschlägigen Abschnitte der Archivfassung
\codeword{paper\_v1\_06\_01\_09\_2025.tex}. Nicht sämtliche mehr als
30.000 Zeilen der aktuellen Theorie samt allen eingeschobenen
Forschungsboxen wurden erneut vollständig geprüft.

Die ursprüngliche Einfachheit ist real auf der Ebene der Konstruktion:
Der Anker \((1,1,2)\) liefert über elementarsymmetrische Funktionen
\((4,5,2)\); der gerade Fockraum von fünf Slots trägt den
16-dimensionalen Halbspinor; D5, A3 und die zyklische Verklebung
organisieren den E8-Rahmen. Die aktuelle Theorie nennt E8 ausdrücklich
eine algebraische Hülle, nicht automatisch die physische Eichgruppe. Die
Quellenmassen und Mischungsgrößen werden über eine kleine Zahl von
Compiler-Ausdrücken organisiert; der physische RG-, QCD- und
kosmologische Transport bleibt davon getrennt.

Die für diese Runde wichtigste wiedergefundene Stelle ist
\codeword{origin\_theory.tex}, Abschnitt zum vierintervalligen freien
Fermionmodell, mit dem Prüfer
\codeword{v480\_multilocal\_four\_interval.py}: Die geometrische
Vierteldrehung hat auf Fermionen die Beziehung

\[\adjustbox{max width=\linewidth}{$\displaystyle \rho^4=-I,\qquad \rho^8=I.$}\]

Auf geraden Paargrößen verschwindet dieses zusätzliche Minus. Der Prüfer
wurde erneut normal und optimiert ausgeführt. Er zeigt den Mechanismus
in einem gewählten antiperiodischen freien Ring; er beweist ausdrücklich
nicht, dass die ursprüngliche TFPT-Seam dieser Ring ist. Sein
Literaturanschluss ist die zustandserhaltende multilokale
Fermionisierung von Rehren und Tedesco:
\href{https://arxiv.org/abs/1205.0324}{Originalarbeit}. Dort ist die
Änderung der Lokalisierung wesentlich; die Identifikation beliebiger
Kopienräume ist nicht frei gegeben.

Die neue Einsicht dieser Runde ist daher nicht, dass Doppelabdeckungen
neu wären. Neu ist der explizite Anschluss dieser bereits bekannten
Information an die überlappenden W-Quellen, den vollständigen
Grundzustandsvergleich und die geladene Antwort.

\subsection{2. Prüfung der zugeschickten
Arbeit}\label{pruxfcfung-der-zugeschickten-arbeit}

Die Arbeit \textbf{TFPT / Universalraum: eine gemeinsame Quelle auf dem
Prüfstand}, Fortsetzung 1.2, enthält 15 Seiten. Ihr Prüfpaket wurde in
einer neuen temporären Kopie entpackt. Alle vier wissenschaftlichen
Prüfer liefen normal und mit \codeword{-OO}; die Ausgaben stimmen
bytegenau mit den gepinnten Ergebnissen überein: 614 Bedingungen,
aufgeteilt in 402, 22, 176 und 14.

Tragende Ergebnisse sind die induzierte Paar-Gram-Matrix, der geschützte
innere Faktor, ein bedingter Wirkungstest, die explizite Beschränkung
der lokalen Kontrollen, Graph-Gegenmodelle, die verdoppelte
Schleifenphase und ein N=2-Grenzwert auf einem vorgegebenen Netz. Die
Reproduktion macht daraus weder 614 unabhängige Theoreme noch eine
vollständige physische Quelle.

Besonders wichtig ist die Gleichheit der gesamten eingeschränkten
bosonischen N=2-Antwort zweier Dreiecksquellen mit Phasen null und pi,
obwohl ihre Fermionkerne verschieden sind. Das neue Paper benennt damit
selbst einen Informationsverlust. Dieser Befund wird im Folgenden nicht
nur wiederholt, sondern für den TFPT-Viererclock verallgemeinert und
durch eine neue Beobachtung ergänzt.

\subsection{3. Ein festgelegter minimaler
Kandidat}\label{ein-festgelegter-minimaler-kandidat}

Vier unabhängige Fermionbereiche besitzen je 64 kanonische Moden
\(f_{x,r}\); vier unabhängige Bosonbereiche besitzen je 60 Moden
\(b_{e,A}\). Die inneren Indizes tragen unverändert die ursprünglichen
Darstellungen \((16,4)\) beziehungsweise den W-Bildraum. Die vier
\textbf{zusätzlichen Bereiche} sind nicht die vier Komponenten der
bereits vorhandenen inneren SU(4)-Darstellung.

Für \(\eta=\pm1\) sei

\[\adjustbox{max width=\linewidth}{$\displaystyle R_\eta=\begin{pmatrix}0&1&0&0\\0&0&1&0\\0&0&0&1\\\eta&0&0&0\end{pmatrix},\qquad
U_\eta=\frac{I+R_\eta}{\sqrt2},\qquad q_{e,r}=\sum_x(U_\eta)_{ex}f_{x,r}.$}\]

Jede einzelne q-Quelle ist kanonisch normiert. Verschiedene q-Quellen
überlappen und sind gerade keine unabhängigen Fermionbanken. Die
gemeinsame Wechselwirkung lautet

\[\adjustbox{max width=\linewidth}{$\displaystyle H_\eta=\Delta N_b+g\sum_{e,A}\bigl(b_{e,A}^\dagger P_A(q_e)+P_A(q_e)^\dagger b_{e,A}\bigr),\qquad
N=N_f+2N_b.$}\]

Dabei werden die 480 ursprünglichen W-Koeffizienten benutzt, keine
zusätzliche freie Sprungmatrix. Der Viereckgraph, die
Zweipunkt-Mittelung, unabhängige Kopien und die Wahl von \(g/\Delta\)
sind zusätzliche, sichtbare Modellannahmen.

\subsubsection{3.1 Der gemeinsame Clock ist eine wirkliche Symmetrie
dieses
Kandidaten}\label{der-gemeinsame-clock-ist-eine-wirkliche-symmetrie-dieses-kandidaten}

\(U_\eta R_\eta=R_\eta U_\eta\). Deshalb transformieren die q-Quellen
unter der Fermiondrehung genauso wie die f-Moden. Beim Paar verschwindet
das Vorzeichen; die Bosonbereiche werden zyklisch ohne Vorzeichen
vertauscht. Die kombinierte Transformation erhält H exakt. Ihre vierte
Potenz ist bei \(\eta=-1\) die gesamte Fermionparität, ihre achte Potenz
die Identität.

Das ist ein expliziter Clock-zu-Quellen-Intertwiner. \textbf{Nicht
gezeigt} ist seine Gleichsetzung mit dem vollständigen mikroskopischen
Seam-Prozess. Insbesondere ist der geometrische C8-Lift nicht durch
Umbenennung der vorhandene innere Clock mit Periode sechs und nicht
automatisch der Coxeter-Clock 30.

\subsubsection{3.2 Paarinformation verliert ein
Vorzeichen}\label{paarinformation-verliert-ein-vorzeichen}

Schreibe \(S=UU^\dagger\). Aus \(WW^\dagger=8I_{60}\) folgt

\[\adjustbox{max width=\linewidth}{$\displaystyle B_{ef}=8S_{ef}^{\,2},\qquad B=8I+2A_{C_4}.$}\]

Es handelt sich um das komplexe Quadrat der Amplitude, nicht generell um
ihr Betragsquadrat. Für die hier reellen Vorzeichen ist B bei beiden
Lifts identisch, mit Eigenwerten \(4,8,8,12\).

Dagegen gilt

\[\adjustbox{max width=\linewidth}{$\displaystyle \operatorname{spec}S_+=\{0,1,1,2\},\qquad
\operatorname{spec}S_-=\{1-1/\sqrt2,1-1/\sqrt2,1+1/\sqrt2,1+1/\sqrt2\}.$}\]

Die periodische Quelle lässt 64 Fermionmoden vollständig frei. H wirkt
auf deren Fockfaktor als Identität; jede globale Energie hat mindestens
die entsprechende \(2^{64}\)-fache Entartung. Das betrifft das gesamte
Fockmodell, nicht nur N=2. Die antiperiodische Quelle hat keinen solchen
linearen Zuschauerbereich.

Allgemeiner hat eine vorgegebene vorzeichenbehaftete Zyklusquelle genau
dann einen eindimensionalen Kopienkern, wenn das Produkt ihrer
Kantenvorzeichen \((-1)^L\) beträgt. Der Beweis ist die Rekursion
\(x_j=-\eta_jx_{j+1}\). Lokale Vorzeichenänderungen verschieben die
Vorzeichen entlang des Zyklus, ändern aber dieses Produkt nicht. Die
zwei Schleifenklassen sind daher keine bloßen verschiedenen Basisnamen.
Exakte Kontrollen umfassen sämtliche Vorzeichenbelegungen der Zyklen mit
drei bis sieben Knoten.

\subsection{4. Die fehlende Information wird ladungsneutral
sichtbar}\label{die-fehlende-information-wird-ladungsneutral-sichtbar}

Die neue Rechnung muss keinen odden Messoperator als frei verfügbaren
Aktuator annehmen. Ein bosonischer Rückkehrzeiger genügt, wenn der
Eingang zusätzlich ein Fermion enthält.

Seien \(A_A\) die reellen antisymmetrischen 64-mal-64-Matrizen zu W. Die
ursprüngliche Dreiladungs-Gram-Matrix gibt

\[\adjustbox{max width=\linewidth}{$\displaystyle K_{Ar,Bs}=(A_BA_A^\dagger)_{rs},\qquad
\operatorname{spec}K=\{8^{(64)},1^{(2880)},(-2)^{(576)},(-4)^{(320)}\}.$}\]

Diese Identität wurde hier erneut aus der vollständigen kanonischen
Dreifermion-Konstruktion und unabhängig aus sämtlichen 3600
A-Matrixblöcken geprüft. Für die zusammengesetzte Quelle definiere

\[\adjustbox{max width=\linewidth}{$\displaystyle T_{(e,x),(f,y)}=S_{ef}\,\overline{U_{fx}}U_{ey}.$}\]

Direktes CAR-Normalordnen liefert auf dem Eingang mit einem Boson und
einem Fermion, nach passender Tensorordnung,

\[\adjustbox{max width=\linewidth}{$\displaystyle \boxed{G_3=(B\otimes I_4)\otimes I_{3840}-T\otimes K.}$}\]

Der 61.440-dimensionale Eingang reduziert sich somit auf vier
16-mal-16-Matrizen \(B\otimes I_4-kT\) mit bekannten Vielfachheiten. Die
Formel wurde zusätzlich mit vollständig ausgeschriebenen kleinen
CAR-Modellen sowie 640 Einträgen des echten 256-Fermion-Modells
kontrolliert; diese Stichprobe wird nicht als vollständige große
Matrixdiagonalisierung bezeichnet.

{\def\LTcaptype{none} % do not increment counter
\begin{longtable}[]{@{}lrr@{}}
\toprule\noalign{}
Moment & Periodisch & Antiperiodisch \\
\midrule\noalign{}
\endhead
\bottomrule\noalign{}
\endlastfoot
\(\operatorname{tr}G_3\) & 487680 & 487680 \\
\(\operatorname{tr}G_3^2\) & 4440960 & 4440960 \\
\(\operatorname{tr}G_3^3\) & 44601120 & 44601120 \\
\(\operatorname{tr}G_3^4\) & 481843440 & 481728240 \\
\end{longtable}
}

Beginnt man im gleichmäßig gemischten Einboson-Einfermion-Eingang und
misst später die gesamte Bosonzahl, folgt

\[\adjustbox{max width=\linewidth}{$\displaystyle \boxed{p_-(t)-p_+(t)=-\frac{g^8t^8}{168}+O(t^{10}).}$}\]

Herleitung: Für einen Eigenwert \(\lambda\) von \(G_3\) ist die
Bosonüberlebenswahrscheinlichkeit

\[\adjustbox{max width=\linewidth}{$\displaystyle 1-\frac{4g^2\lambda}{\Delta^2+4g^2\lambda}
\sin^2\!\left(\frac{t\sqrt{\Delta^2+4g^2\lambda}}2\right).$}\]

Der oberste \(\lambda^4\)-Term der achten Zeitordnung ist
\(g^8\lambda^4/315\). Die niedrigeren Momente stimmen überein; Division
der Differenz \(-115200\) durch \(315\cdot61440\) ergibt \(-1/168\).
Dies ist eine exakte lokale Tayloraussage, keine numerisch zertifizierte
Differenz bei beliebigen endlichen Zeiten.

\textbf{Wichtig:} Dieser N=3-Vergleich ist ein
Quellen-Diagnoseexperiment auf einem gesetzten Eingang. Er ist nicht die
Antwort auf dem folgenden N=256-Grundzustand. Diese beiden Aufgaben
werden ausdrücklich getrennt.

\subsection{5. Eine vollständige lokale Casimir-Identität ersetzt
riesige
Fockmatrizen}\label{eine-vollstuxe4ndige-lokale-casimir-identituxe4t-ersetzt-riesige-fockmatrizen}

Für eine kanonische 64-Moden-Quelle sei

\[\adjustbox{max width=\linewidth}{$\displaystyle Q=\sum_A P_A^\dagger P_A.$}\]

Die 60 antihermiteschen Einteilchengeneratoren \(X_a\) von Spin(10) mal
SU(4) erfüllen \(-\sum_aX_a^2=60I\). Auf dem gesamten Fockraum gilt

\[\adjustbox{max width=\linewidth}{$\displaystyle \boxed{8Q=60N_f-\mathcal C,\qquad
\mathcal C=\sum_a d\Gamma(X_a)^\dagger d\Gamma(X_a)\ge0.}$}\]

Warum die endliche Prüfung für diese Identität reicht: Nach
CAR-Normalordnung haben beide Seiten höchstens Grad vier. Die
Einteilchenkoeffizienten sind durch den vollständigen Einteilchencasimir
festgelegt. Die quartischen Koeffizienten werden durch die vollständige
Zweifermionmatrix bestimmt. Dort wurde exakt

\[\adjustbox{max width=\linewidth}{$\displaystyle 8W^\dagger W+\mathcal C_{\Lambda^2}-120I_{2016}=0$}\]

nachgerechnet. Vakuum, Einteilchen- und Zweiteilchenteil bestimmen daher
alle normalgeordneten Koeffizienten; der Satz gilt auf allen \(2^{64}\)
Fermionzuständen. Die Matrixrechnung verwendet exakte gaußsche
Ganzzahlen ohne Toleranzvergleich. Dies ist eine erneute Verifikation
der hier benötigten Casimir-Grundlage, nicht eine Übernahme des alten
Grundzustandssatzes für die einzelne Bank.

Daraus folgt \(Q\le(15/2)N_f\le480I\). Der voll besetzte Zustand ist ein
Singulett und erreicht 480. Für vier überlappende Quellen gilt deshalb

\[\adjustbox{max width=\linewidth}{$\displaystyle Q_{\rm ges}=\sum_{e,A}P_A(q_e)^\dagger P_A(q_e)\le M I,\qquad M=1920.$}\]

Mit \(s_0=1-1/\sqrt2\) und \(\delta=(15/2)s_0>2\) erhält man für die
antiperiodische Quelle die stärkere operatorielle Aussage

\[\adjustbox{max width=\linewidth}{$\displaystyle \boxed{M-Q_{\rm ges}\ge\delta N_h,\qquad N_h=256-N_f.}$}\]

Denn \(480-Q_e\ge(15/2)(64-N_{q_e})\), und die Summe der lokalen
Lochzahlen ist die zweite Quantisierung von \(U^\dagger U\) im Lochraum;
deren kleinster Eigenwert ist \(s_0\). Der vollständig gefüllte globale
Zustand \(F\) erreicht M und ist der einzige Maximierer. Das ist bereits
die eindeutige Zustandswahl der effektiven Wechselwirkung
\(-(g^2/\Delta)Q_{\rm ges}\).

\subsection{6. Vollständiger Grundzustandssatz ohne
Bosonabschneiden}\label{vollstuxe4ndiger-grundzustandssatz-ohne-bosonabschneiden}

\textbf{Satz.} Für die festgelegte antiperiodische Viererquelle,
\(\Delta>0\), \(\mu=0\) und

\[\adjustbox{max width=\linewidth}{$\displaystyle 0<|g|/\Delta\le1/2000$}\]

hat H auf dem vollständigen Fermion-Boson-Fockraum genau einen
Grundzustand. Er liegt in der erhaltenen Gesamtladung \(N=256\). Seine
spektrale Lücke ist größer als \(g^2/\Delta\).

\subsubsection{6.1 Die drei einfachen
Schranken}\label{die-drei-einfachen-schranken}

Quadratvervollständigung ergibt für beliebiges \(0<\theta\le1\)

\[\adjustbox{max width=\linewidth}{$\displaystyle H\ge(1-\theta)\Delta N_b-\frac{g^2}{\theta\Delta}Q_{\rm ges}.$}\]

\textbf{Kleinere Gesamtladungen.} Für \(N\le255\) gilt \(N_h\ge1\), also
bei \(\theta=1\)

\[\adjustbox{max width=\linewidth}{$\displaystyle H|_{N\le255}\ge-(M-\delta)g^2/\Delta.$}\]

\textbf{Größere Gesamtladungen.} Für \(N\ge257\) muss mindestens ein
Boson vorhanden sein. Mit \(\theta=1/2\) folgt

\[\adjustbox{max width=\linewidth}{$\displaystyle H|_{N\ge257}\ge\Delta/2-2Mg^2/\Delta>0.$}\]

\textbf{Gesamtladung 256.} Der bosonfreie Teil dieses Sektors ist
eindimensional, nämlich F. Auf seinem orthogonalen Komplement gilt
dieselbe positive Schranke wie oben. Auf dem Zweizustandsraum aus F und
seiner normierten ersten Umwandlung ist H genau die Ritz-Matrix

\[\adjustbox{max width=\linewidth}{$\displaystyle \begin{pmatrix}0&g\sqrt M\\g\sqrt M&\Delta\end{pmatrix}.$}\]

Sie liefert die obere Grundzustandsschranke

\[\adjustbox{max width=\linewidth}{$\displaystyle E_{\rm Ritz}=\frac{\Delta-\sqrt{\Delta^2+4Mg^2}}2
\le-M\frac{g^2}{\Delta}+M^2\frac{g^4}{\Delta^3}<0.$}\]

Die Trennung von den kleineren Ladungen ist mindestens

\[\adjustbox{max width=\linewidth}{$\displaystyle \left(\delta-M^2\frac{g^2}{\Delta^2}\right)\frac{g^2}{\Delta}>
\frac{g^2}{\Delta},$}\]

weil \(M^2/2000^2=0.9216\) und \(\delta>2\). Innerhalb N=256 kann es
höchstens einen negativen Eigenwert geben: Das Komplement der einen
bosonfreien Richtung ist positiv. Das Min-Max-Prinzip liefert die
Eindeutigkeit und die behauptete Lücke.

\subsubsection{6.2 Warum dabei keine unendlichen Zustände unterschlagen
werden}\label{warum-dabei-keine-unendlichen-zustuxe4nde-unterschlagen-werden}

Es gibt endlich viele Fermion- und Bosonmoden, aber beliebige
Bosonbesetzungen. Die Kopplung ist relativ zu \(N_b\) infinitesimal
beschränkt: Auf dem Übergang von k zu k+1 Bosonen hat
\(\sum b^\dagger P\) die Normschranke \(\sqrt{M(k+1)}\). Daraus folgen
Selbstadjungiertheit auf der Zahloperator-Domäne und kompakte
Resolvente. Die oberen Ladungssektoren sind durch die explizite Schranke
gleichzeitig ausgeschlossen, nicht durch eine endliche Stichprobe. Alle
Aussagen betreffen den vollständigen H dieses Kandidaten.

Dieser Beweis ist wesentlich kleiner als eine Vollraumberechnung. Er
benutzt die eigentliche TFPT-Stärke - Darstellung, Gradierung und
Erhaltung - um die große Matrix gar nicht erst aufzustellen.

